

\Bsec{Lineare Gleichungssysteme}{
		
		\begin{itemize}
			\itemf Gauß-Elimination
			\itemf LR-Zerlegung
			\itemf QR-Zerlegung
			\itemf Lineares Ausgleichsproblem
			\itemf Kombinations Aufgaben und Verständnisfragen
		\end{itemize}% \sbreak
%		Noch nicht vorbereitet auf dieses Kapitel? Dann 
		}% {}
		
%----------------------------------------------------------------------------------
\bsubsec{Gauß-Elimination}{
%		
%		Man möchte ein Problem der Form $Ax = b$ lösen.
%		
%	}
%		
		Löse folgende LGS mithilfe von Gauß-Elimination:
		
		\begin{enumerate}
			\itemf 
			$\begin{bmatrix}
				3 & 0 & 1 & 0 & | & 1\\
				0 & 3 & 2 & 1 & | & 4\\
				1 & 0 & 1 & 0 & | & 2\\
				0 & 1 & 0 & 1 & | & 2
			\end{bmatrix}$
			\itemf 
			$\begin{bmatrix}
				1 & 2 & | & 2\\
				a & 4 & | & b\\
				3 & 5 & | & 5
			\end{bmatrix}$ ($a$ und $b$ sind Variablen)
			\itemf
			$\begin{bmatrix}
				2 & 5 & 2 & | & 1\\
				1 & 4 & 2 & | & 2
			\end{bmatrix}$
		\end{enumerate}
		
	}
	\bsol{
		
		\begin{enumerate}
			\itemf $x =(-0.5; -1.5; 2.5; 3.5)\trans$
			\itemf Keine Lösung
			\itemf Überbestimmtes LGS $\Rightarrow$ Für uns so nicht lösbar!
		\end{enumerate}
	}
%----------------------------------------------------------------------------------
\bsubsec{LR-Zerlegung}{
%		
%		$$\begin{bmatrix}
%		& & \\
%		& A& \\
%		& & 
%		1 & 0 & 0\\
%		& 1 & 0\\
%		&   & 1
%		&&&\\
%		0&&\\
%		0&0&
%		$$
%	}
%		
		Löse folgende LGS mithilfe von LR-Zerlegung:
		
		\begin{enumerate}
			\itemf $Ax = b$ mit $A = 
			\begin{bmatrix}
			2 & 0 & 1 & 0\\
			0 & 2 & 2 & 1\\
			1 & 0 & 1 & 0\\
			0 & 1 & 0 & 1
			\end{bmatrix};\quad b = 
			\begin{bmatrix}
				1 \\ 4 \\ 2 \\ 2
			\end{bmatrix}$
			\itemf $Az = c$ mit $c =
			\begin{bmatrix}
				0 \\ 2 \\ 4 \\ 4
			\end{bmatrix}$
		\end{enumerate}
		
	}
	\bsol{
		
		\begin{enumerate}
			\itemf $x = (-1; -4; 3; 6)\trans$ mit L = 
			$\begin{bmatrix}
				1 & 0 & 0 & 0 \\
				0 & 1 & 0 & 0 \\
				0.5 & 0 & 1 & 0 \\
				0 & 0.5 & -2 & 1
			\end{bmatrix}$ und R = 
			$\begin{bmatrix}
				2 & 0 & 1 & 0 \\
				0 & 2 & 2 & 1 \\
				0 & 0 & 0.5 & 0 \\
				0 & 0 & 0 & 0.5
			\end{bmatrix}$
			\itemf $z = (-4; -18; 8; 22)\trans$ mit obigen L und R
		\end{enumerate}
	}

%----------------------------------------------------------------------------------
\bsubsec{QR-Zerlegung}{
%		
%		Mit $(a,b)^T$ als die erste Spalte der 2x2 Matrix A, $$c = \frac{a}{\sqrt{a^2 + b^2}}, s = \frac{b}{\sqrt{a^2 + b^2}}$$
%		$$G_\varphi = \begin{bmatrix}
%		c & s\\ -s & c
%		$$\begin{bmatrix}
%		c & s\\ -s & c
%		a\\b
%		r\\ 0
%		$$A = QR \Leftrightarrow Q^{-1} A = R \Leftrightarrow G_\varphi A = R$$
%		$$Rx = Q^T \cdot b$$
%	}
%		
		Zerlege die Matrix A mit der QR-Zerlegung und rechne das Ergebnis des LGS mithilfe der Zerlegung aus:
		
		\begin{enumerate}
			\itemf $Ax = b$ mit $A = 
			\begin{bmatrix}
				3 & 2 \\
				4 & 1
			\end{bmatrix}; \quad b = 
			\begin{bmatrix}
				5 \\ 6
			\end{bmatrix}$
			\itemf $Ax = b$ mit $A = 
			\begin{bmatrix}
			0 & 2 & 2\\
			1 & 1 & 1\\
			1 & 0 & 2
			\end{bmatrix}; \quad b = 
			\begin{bmatrix}
			1 \\ 2 \\ 3
			\end{bmatrix}$
		\end{enumerate}
		
	}
	\bsol{
		
		\begin{enumerate}
			\item $x = (\frac{7}{5}; \frac{2}{5})\trans$ mit 
				$$Q = 
				\begin{bmatrix}
					\frac{3}{5} & -\frac{4}{5} \\
					\frac{4}{5} & \frac{3}{5}
				\end{bmatrix}; \quad R = 
				\begin{bmatrix}
					5 & 2 \\
					0 & -1
				\end{bmatrix}$$
			\item $x=(0; 0.5; 1.5)\trans$ mit $$Q =
			\begin{bmatrix}
				0 					& \frac{4}{3 \sqrt{2}} & -\frac{1}{3} \\
				\frac{1}{\sqrt{2}}	& \frac{1}{3 \sqrt{2}} & \frac{2}{3} \\
				\frac{1}{\sqrt{2}}	& -\frac{1}{3 \sqrt{2}}& -\frac{2}{3}
			\end{bmatrix}; \quad R =
			\begin{bmatrix}
				\sqrt{2}	& \frac{1}{\sqrt{2}}	& \frac{3}{\sqrt{2}} \\
				0			& -\frac{1}{3 \sqrt{2}} & -\frac{7}{3 \sqrt{2}} \\
				0 & 0 		& -\frac{4}{3}
			\end{bmatrix}$$
		\end{enumerate}
	}
%----------------------------------------------------------------------------------
\bsubsec{Lineares Ausgleichsproblem}{
%	Allgemeine Geradengleichung:
%	$$y = m \cdot x + t$$
%	Die Normalengleichung: $$ A^T \cdot b = A^T \cdot A \cdot x' $$
%	Mit QR Zerlegung:
%	$$Rx = \beta_1$$ 
%	$$\begin{bmatrix}
%	Die rechte Seite der Gleichung entspricht \fat{NICHT} einem Bruch, sondern der Aufteilung des Vektores in einen oberen ($=\beta_1$) und unteren ($=\beta_2$) Teil.
%	
%}
%	
	Es soll eine Gerade bestimmt werden, die möglichst gut durch folgende Punkte geht:
	$$P_0(1|-3) \aae P_1(2|3) \aae P_2(4|-2)$$
	
	\begin{enumerate}
		\itemf Stelle das LGS zu obigen Problem auf.
		\itemf Löse obiges Problem mit der Normalengleichung.
	\end{enumerate}
	
}
\bsol{
	
	\begin{enumerate}
		\itemf $$\begin{bmatrix}
		1 & 1 \\
		1 & 2 \\
		1 & 4 \\
		\end{bmatrix} \cdot \begin{bmatrix}
		t \\ m
		\end{bmatrix} = 
		\begin{bmatrix}
		-3 \\ 3 \\ -2
		\end{bmatrix}$$
		\itemf $$\begin{bmatrix}
			3 & 7 &|& -2 \\
			7 & 21 &|& -5
		\end{bmatrix} \longrightarrow \begin{bmatrix}
		3 & 7 &|& -2 \\
		-2 & 0 &|& 1
		\end{bmatrix}$$
		$$t = -\frac{1}{2} \aae m = -\frac{1}{14}$$
		$$y = -\frac{1}{14}x - \frac{1}{2}$$
	\end{enumerate}
}
%\btask{
%	
%	\begin{enumerate}
%		\item Löse folgendes Lineares Ausgleichsproblem mit der QR-Zerlegung.
%		$$\begin{bmatrix}
%		1 & 1 \\
%		1 & 2 \\
%		1 & 0 \\
%		\end{bmatrix} \cdot \begin{bmatrix}
%		t \\ m
%		\end{bmatrix} = 
%		\begin{bmatrix}
%		4 \\ 2 \\ 1
%		\end{bmatrix}$$
%	\end{enumerate}
%	
%}
%\bsol{
%	
%	\begin{enumerate}
%		\item 
%	\end{enumerate}
%}
%----------------------------------------------------------------------------------
\bquestion{
%	}
%		
		Beantworte folgende (Trick-)Fragen möglichst genau:
		\begin{enumerate}
			\itemf Was ist \fat{Pivotisierung} und wie funktioniert es?
			\itemf Gauß, LR und QR: Was ist am stabilsten? Warum?
			\itemf Wenn du mit derselben Matrix ganz viele LGS lösen musst, welches der drei Verfahren wendest du an, wenn der Fokus auf der Laufzeit liegt?
			\itemf Du hast eine Matrix der Form: $\begin{bmatrix}
				a & b & c\\
				0 & d & e\\
				0 & 0 & f
			\end{bmatrix}$ \sbreak
			Du sollst die Matrix in Q und R zerlegen, was fällt dir auf?
		\end{enumerate}
		
	}
	\bsol{
		
		\begin{enumerate}
			\itemf siehe mein Skript \Smiley
			\itemf QR-Zerlegung, weil dort mit orthogonalen Matrizen gerechnet wird.
			\itemf LR-Zerlegung: QR ist zwar stabiler, aber leicht kostenaufwändiger.
			\itemf Wir brauchen Givens-Rotation nicht anzuwenden, weil sich die Matrix schon in einer rechts oberen Dreiecksmatrix befindet (Sprich man kann direkt mit Rückwärtssubst. anfangen).
		\end{enumerate}
	}
%----------------------------------------------------------------------------------